% ============================================================
% SS-4: String Tension from the 600-Cell Face-Mode Multiplicity
%       (Resolving OPEN-SS-5: "String Tension sigma from sea_strength")
% Version 0.1 — 16 April 2026
%
% CHANGELOG:
% v0.1 (16 Apr 2026): Initial draft. Presents the conjecture
%   sigma = M_0 * z^2 / (phi * l_edge), which corrects the
%   heuristic factor in SS-2 (z*pi -> z^2) on the physical grounds
%   that the ZBW orbit is a discrete z-polygon on the lattice,
%   not a continuous circle. Numerical agreement with the
%   charmonium Cornell fit: +1.82%. Cornell closure consistency
%   for r_conf: -1.6%. The rigorous derivation of the z^2
%   face-mode multiplicity from the 600-cell trace structure
%   (Tr A^2, Tr A^3) is deferred to Section 10 as the remaining
%   open step. Includes Physical Mechanism Bridge per new SS-3
%   template. Problem ID: this paper moves OPEN-SS-5 to CONJ-SS-5
%   (problem number 5 retained through the progression per CPP
%   nomenclature). Paper ID: SS-4, the next sequential paper in
%   the Strong Sector series (paper numbering is independent of
%   problem numbering). Registers OPEN-SS-15 (rigorous z^2
%   multiplicity) and marks CONJ-SS-2-1 as superseded.
%   Added Remark 8.2 (4+4 standing-wave extension), flagging
%   the natural next question: baryon size from 8-mode resonance
%   with sigma as tension input.
%
% Series: Strong Sector
% Dependencies: SS-1 (sea_strength, gauge trace structure),
%               SS-2 (l_unit, l_edge, sigma heuristic),
%               SS-3 (SU(3) uniqueness, face-mode basis),
%               SM-8 (M_0 = m_e z/phi, DP energy quantum),
%               C14 (Cornell closure, self-consistency)
% Resolves: OPEN-SS-5 (modulo rigorous step in Section 10)
% Repository: CPP/series_strong/papers/SS-4_string_tension.tex
% ============================================================

\documentclass[11pt,letterpaper]{article}
\usepackage[utf8]{inputenc}
\usepackage[margin=1in]{geometry}
\usepackage[T1]{fontenc}
\usepackage{amsmath,amssymb,amsfonts,amsthm}
\usepackage{graphicx}
\usepackage{xcolor}
\usepackage{hyperref}
\usepackage{booktabs}
\usepackage{array}
\usepackage{enumitem}
\usepackage{float}
\usepackage[font=small, labelfont=bf]{caption}
\usepackage{parskip}
\usepackage{tikz}

% Bibliography
\usepackage[authoryear,round]{natbib}
\bibliographystyle{plainnat}

% Hyperlinks
\hypersetup{
    colorlinks=true,
    linkcolor=blue,
    citecolor=blue,
    urlcolor=blue
}

% Theorem environments
\newtheorem{theorem}{Theorem}[section]
\newtheorem{proposition}[theorem]{Proposition}
\newtheorem{lemma}[theorem]{Lemma}
\newtheorem{corollary}[theorem]{Corollary}
\newtheorem{conjecture}[theorem]{Conjecture}
\newtheorem{definition}[theorem]{Definition}
\newtheorem{remark}[theorem]{Remark}
\newtheorem*{openstep}{Open Step}

% Standard CPP macros
\newcommand{\lP}{l_P}
\newcommand{\tP}{t_P}
\newcommand{\EP}{E_P}
\newcommand{\SSV}{\mathrm{SSV}}
\newcommand{\phig}{\varphi}
\newcommand{\Tr}{\operatorname{Tr}}

% Paper-specific macros
\newcommand{\ledge}{l_{\mathrm{edge}}}
\newcommand{\lunit}{l_{\mathrm{unit}}}
\newcommand{\LQCD}{\Lambda_{\mathrm{QCD}}}
\newcommand{\alphageom}{\alpha_{\mathrm{geom}}}
\newcommand{\kSM}{k_{\mathrm{SM}}}
\newcommand{\sea}{\mathrm{sea\_strength}}
\newcommand{\rconf}{r_{\mathrm{conf}}}
\newcommand{\GeVfm}{\mathrm{GeV/fm}}
\newcommand{\MeVfm}{\mathrm{MeV/fm}}

% ============================================================

\title{\textbf{SS-4: String Tension from the 600-Cell Face-Mode Multiplicity}\\[6pt]
{\large Conscious Point Physics --- Strong Sector Series}\\[4pt]
{\normalsize Version 0.1, 16 April 2026 \ (DRAFT)}}

\author{Thomas Lee Abshier, ND \and Claude Opus (Anthropic)}

\date{\normalsize Hyperphysics Institute\\
\url{https://hyperphysics.com}\\
\href{mailto:drthomas007@protonmail.com}{drthomas007@protonmail.com}}

\begin{document}
\maketitle

% ===================================================================
\begin{abstract}
SS-2 introduced a heuristic formula for the QCD string tension,
$\sigma_{\mathrm{SS\text{-}2}} = M_0\,z\pi / (\phig\,\ledge)
\approx 243~\MeVfm$, justifying the factor of $\pi$ by the
``ZBW circular orbit'' but obtaining a value $\sim 3.8\times$
too small relative to the charmonium Cornell fit
$\sigma \approx 910~\MeVfm$.  We identify the source of the
discrepancy: the ZBW orbit in CPP is not a continuous circle but
a \emph{discrete lattice polygon} traversing $z = 12$
nearest-neighbour vertices on the 600-cell.  The continuum
circumference factor $2\pi$ is replaced by the discrete vertex
count $z$, giving one power of $z$; and the mode multiplicity
around the closed orbit contributes a second power, producing
the factor $z^2$.  The resulting conjecture
\begin{equation*}
\boxed{\sigma = \frac{M_0 \cdot z^2}{\phig \cdot \ledge} = 926.5~\MeVfm}
\end{equation*}
agrees with the charmonium Cornell fit to $+1.82\%$, and the
self-consistent confinement radius $\rconf = \sqrt{\alpha_s \hbar c /
\sigma} = 0.158~\mathrm{fm}$ agrees with the C14 closure value to
$-1.6\%$.  Both emerge simultaneously with no parameter fitting
beyond the CPP primitives already present in SS-1 and SS-2
($m_e$, $z$, $\phig$, $\alphageom$, from which $\sea$,
$\LQCD$, $\lunit$, $\ledge$, and $M_0$ all follow).

The factorisation $\sigma = (M_0/\ledge) \times (z^2/\phig)$
admits a clean physical reading: the flux tube carries
one DP energy quantum $M_0$ per lattice edge of length
$\ledge$, with a face-mode multiplicity of $z^2$ (coordination
along the tube axis times coordination around the tube
cross-section) and a golden-ratio propagation efficiency
$\eta = 1/\phig$ inherited from the 600-cell edge-to-circumradius
ratio.  The rigorous derivation of the $z^2$ multiplicity from
the 600-cell gauge traces $\Tr A^2 = 1440$ and $\Tr A^3 = 7200$
is stated as the remaining open step (Section~\ref{sec:open}).

\medskip\noindent
\textbf{Open Problem resolved (modulo open step):} OPEN-SS-5
(string tension $\sigma$ from $\sea$).

\medskip\noindent
\textbf{Consequence:} SS-2's heuristic
$\sigma = M_0 z\pi/(\phig \ledge)$ is superseded by
$\sigma = M_0 z^2/(\phig \ledge)$, the $z/\pi \approx 3.82$
correction factor is physically identified, and the SS-2 proton
radius derivation inherits a quantitative string tension at the
2\% level.
\end{abstract}

\textbf{Keywords:} QCD string tension, Cornell potential, flux tube,
confinement radius, qDP chain, 600-cell polytope, face-mode multiplicity,
ZBW (zitterbewegung) orbit, discrete gauge theory, Conscious Point
Physics, sea\_strength, bow rigidity

\textbf{Plain Language Summary:}
Quarks inside protons and mesons are held together by a force
that grows with distance --- the ``string tension'' of roughly
$0.9~\GeVfm$, meaning it takes about one gigaelectronvolt of
energy to pull two quarks one femtometre apart.  Where does
this specific number come from?  In standard QCD it is measured
experimentally and fitted into the Cornell potential.  In CPP,
the quark binding is a physical chain of Dipole Pair (DP) bonds
laid down on a regular lattice.  An earlier CPP paper (SS-2)
proposed that the string tension equals the DP energy quantum
times the number of orbital steps times $\pi$, divided by the
golden ratio and the lattice edge length --- a formula that came
out four times too small.  This paper identifies the error:
the orbit inside a quark is not a continuous circle but a
discrete 12-vertex polygon on the lattice, so the continuum
$\pi$ should be replaced by the vertex count $z = 12$, and the
orbit contributes $z^2$ rather than $z\pi$.  With that single
correction the formula gives $0.93~\GeVfm$, within 2\% of
the experimentally fitted value, and the characteristic
confinement distance of 0.16 femtometres falls out as a bonus.

\raggedright
\tableofcontents
\newpage

% ===================================================================
\section{Introduction}
\label{sec:intro}
% ===================================================================

The Cornell potential
\begin{equation}
V(r) = -\frac{\alpha_s \hbar c}{r} + \sigma\, r,
\qquad
\sigma \approx 0.9~\GeVfm
\label{eq:cornell}
\end{equation}
provides the phenomenological account of quark confinement in
heavy quarkonium and, through lattice QCD, the asymptotic linear
behaviour of the quark-antiquark potential at intermediate
distances~\citep{eichten_quigg, bali_review}.  The coefficient
$\sigma$, called the QCD string tension, is a fundamental
dimensional quantity of the strong sector but in both the Cornell
model and standard QCD it is taken from experiment or lattice
simulation rather than derived.

In Conscious Point Physics~\citep{cpp_ss1}, confinement is the
energetic stability of a self-collimated qDP chain laid down
between colour-separated quarks on the 600-cell lattice.  The
qualitative mechanism --- bow rigidity of the alternating
compressive/tensile DP chain --- was established in the
\texttt{chain\_fraying\_dynamics} notebook (January 2026;
Stage~13 of the Strong Sector development record).  The
self-consistent closure
\begin{equation}
\sigma = \frac{\alpha_s \hbar c}{\rconf^2},
\qquad
\rconf = \sqrt{\frac{\alpha_s \hbar c}{\sigma}},
\label{eq:cornell_closure}
\end{equation}
(companion paper C14~\citep{c14}) fixes $(\sigma, \rconf)$ given
$\alpha_s$, but does not deliver $\sigma$ from CPP primitives
alone: it is one equation in two unknowns.

\subsection{The puzzle in SS-2}
\label{ssec:puzzle}

SS-2~\citep{cpp_ss2_lattice} introduced the heuristic
\begin{equation}
\sigma_{\mathrm{SS\text{-}2}}
= \frac{M_0 \cdot z\pi}{\phig \cdot \ledge}
= \frac{3.790~\mathrm{MeV} \times 12\pi}{1.618 \times 0.364~\mathrm{fm}}
= 243~\MeVfm,
\label{eq:ss2_heuristic}
\end{equation}
with $M_0 = m_e z/\phig = 3.790~\mathrm{MeV}$ (the DP energy
quantum~\citep{cpp_sm8}), $z = 12$ (coordination number of the
600-cell), $\phig = (1+\sqrt{5})/2$ (golden ratio), and
$\ledge = \lunit/\phig = 0.364~\mathrm{fm}$ (lattice edge
length, derived in SS-2 from $\LQCD$).  The factor of $\pi$ was
motivated by the ``ZBW circular orbit'' of the qCP charges in
the chain.  SS-2 explicitly flagged the formula as
``physically motivated but not rigorously derived,'' with the
remark that the specific decomposition of $z\pi/\phig$ required
justification and that deriving $\sigma$ from the 600-cell
mode spectrum was an open problem.

Numerically, the heuristic sits a factor of
$\sigma_{\mathrm{Cornell}}/\sigma_{\mathrm{SS\text{-}2}}
\approx 910/243 \approx 3.75$ below the target.  The size of
this discrepancy is suggestive: the ratio $z/\pi = 12/\pi
\approx 3.820$ is within $1.9\%$ of the missing factor,
hinting that a single symbol was wrong in
\eqref{eq:ss2_heuristic}.

\subsection{This paper}

We identify the error.  The ``ZBW circular orbit'' in
\eqref{eq:ss2_heuristic} is not a continuous circle but a
discrete lattice polygon with $z = 12$ vertices; the continuum
circumference factor $2\pi$ has its discrete analogue in the
vertex count $z$, and the mode multiplicity around a closed
polygon orbit contributes a second factor of $z$.  Replacing
$z\pi$ with $z^2$ gives
\begin{equation}
\sigma = \frac{M_0 \cdot z^2}{\phig \cdot \ledge}
= 926.5~\MeVfm,
\label{eq:main}
\end{equation}
within $+1.82\%$ of the charmonium Cornell fit $0.910~\GeVfm$
and self-consistent with the C14 closure:
$\rconf = \sqrt{\alpha_s \hbar c / \sigma} = 0.158~\mathrm{fm}$
at $\alpha_s = 0.118$, vs.\ the C14 value $0.161~\mathrm{fm}$
(error $-1.6\%$).

The numerical match is presented alongside three cautions.
First, the replacement of $\pi$ by $z$ is physically motivated
by the discrete-vs-continuum distinction; the rigorous
derivation of the $z^2$ multiplicity from the 600-cell gauge
traces $\Tr A^2 = 1440$ and $\Tr A^3 = 7200$ is deferred to
Section~\ref{sec:open} as the remaining open step.  Until that
step is completed, \eqref{eq:main} is registered as a
Conjecture (Section~\ref{sec:main}) rather than a Theorem.
Second, $\sea$ enters $\sigma$ indirectly through $\LQCD$
(and thence $\lunit$, $\ledge$), not as a direct multiplicative
factor; the phrasing of OPEN-SS-5 (``$\sigma$ from
$\sea$'') should be read in this extended sense.  Third, the
bow-rigidity mechanism established in the
\texttt{chain\_fraying\_dynamics} notebook is a qualitative
confirmation that $V(r)$ is linear in $r$; the \emph{magnitude}
of $\sigma$ is set by the mode-count formula, not by the bow
amplitude.

The paper is organised as follows.  Section~\ref{sec:primitives}
lists the CPP primitives that enter $\sigma$.
Section~\ref{sec:puzzle} documents SS-2's heuristic and the
$z/\pi$ puzzle.  Section~\ref{sec:resolution} presents the
discrete-polygon resolution.  Section~\ref{sec:main} states the
conjecture and verifies the numerical match.
Section~\ref{sec:factorization} gives the physical factorisation
$\sigma = (M_0/\ledge)(z^2/\phig)$.  Section~\ref{sec:closure}
checks Cornell-closure consistency.  Section~\ref{sec:bridge}
presents the Physical Mechanism Bridge from CPP to QCD.
Section~\ref{sec:predictions} lists falsifiable predictions
inherited from the \texttt{chain\_fraying} and
\texttt{zbw\_magnetic} notebooks.  Section~\ref{sec:open} states
the remaining rigorous step.  Appendix~A documents the numerical
verification.

% ===================================================================
\section{CPP Primitives and Notation}
\label{sec:primitives}
% ===================================================================

We take as given:
\begin{itemize}[leftmargin=1.5em]
\item $m_e = 0.511~\mathrm{MeV}$ (electron mass, the sole mass
  calibration for the Standard Model emergence
  series~\citep{cpp_sm8}).
\item $z = 12$ (coordination number of the 600-cell polytope:
  each vertex has exactly 12 nearest neighbours~\citep{cpp_ss1}).
\item $\phig = (1+\sqrt{5})/2 \approx 1.6180$ (golden ratio;
  ratio of 600-cell circumradius to edge length, and
  propagation-efficiency constant $\eta = 1/\phig$).
\item $\alphageom = 3(11 + 5\sqrt{5})\sqrt{5+\sqrt{5}}/320
  \approx 0.5594$ (600-cell Voronoi stiffness integral,
  Theorem~11.1 of SS-1).
\item $\LQCD = 335~\mathrm{MeV}$ (derived in SS-2 via the
  $f_\pi$ route $\LQCD = 2\pi f_\pi/\sqrt{3}$ and confirmed
  via $\alphageom = 1/\sqrt{5}$ running, both converging to
  the same value).
\item $\hbar c = 0.1973~\mathrm{GeV}\cdot\mathrm{fm}$ (conversion
  constant).
\end{itemize}
From these the following derived quantities appear throughout:
\begin{align}
\kSM     &= \frac{\alphageom}{z\,\phig^2} \approx 0.01780,
  \label{eq:kSM}\\
\sea     &= \frac{N_{\mathrm{lattice}}}{z} \cdot \kSM
          = 10\,\kSM \approx 0.1780,
  \label{eq:sea_strength}\\
M_0      &= \frac{m_e\, z}{\phig} = 3.790~\mathrm{MeV},
  \label{eq:M0}\\
\lunit   &= \frac{\hbar c}{\LQCD} = 0.589~\mathrm{fm},
  \label{eq:lunit}\\
\ledge   &= \frac{\lunit}{\phig} = 0.364~\mathrm{fm},
  \label{eq:ledge}
\end{align}
where $N_{\mathrm{lattice}} = 120$ is the 600-cell vertex count
and the factor of $10$ in \eqref{eq:sea_strength} is the
lattice-shell multiplicity $N_{\mathrm{lattice}}/z$ (SS-1,
Theorem~11.2).  All six quantities trace back to
$\{m_e, z, \phig, \alphageom\}$ and hence contain no
new calibration beyond the electron mass.

% ===================================================================
\section{SS-2's Heuristic and the \texorpdfstring{$z/\pi$}{z/pi} Puzzle}
\label{sec:puzzle}
% ===================================================================

SS-2 Section~4 proposed
\begin{equation}
\sigma_{\mathrm{SS\text{-}2}}
= \frac{M_0 \cdot z\pi}{\phig \cdot \ledge}
= 243~\MeVfm,
\label{eq:heuristic}
\end{equation}
with the narrative that ``$\pi$ enters from the ZBW circular
orbit''~\citep{cpp_ss2_lattice}.  The author noted that the
formula used a specific decomposition of $z\pi/\phig$ that was
``physically motivated but not rigorously derived,'' and listed
``derive $\sigma$ rigorously from the 600-cell lattice mode
spectrum, eliminating the $z\pi/\phig$ ansatz'' in the
open-problems table.

The puzzle: the charmonium Cornell fit gives
$\sigma \approx 910~\MeVfm$, a factor
\begin{equation}
R := \frac{\sigma_{\mathrm{Cornell}}}{\sigma_{\mathrm{SS\text{-}2}}}
   = \frac{910}{243} \approx 3.745
\end{equation}
larger.  Across candidate symbolic factors, $R$ is closest to
$z/\pi = 12/\pi \approx 3.820$:
\begin{center}
\begin{tabular}{lc}
\toprule
Factor & Numerical value \\
\midrule
$R = \sigma_{\mathrm{Cornell}}/\sigma_{\mathrm{SS\text{-}2}}$
  & $3.745$ \\
$z/\pi$ & $3.820$ \\
$\phig^3 = \phig^2 + \phig$ & $4.236$ \\
$4/\phig^{1/2}$ & $3.145$ \\
$\phig\cdot e$ & $4.399$ \\
\bottomrule
\end{tabular}
\end{center}
The $z/\pi$ match is within $2.0\%$, stronger than any
competitor by a substantial margin.  This is the suggestive
clue that $\sigma_{\mathrm{SS\text{-}2}}$ is missing a factor
of $z/\pi$ --- equivalently, that $z\pi$ in \eqref{eq:heuristic}
should have been $z^2$.

% ===================================================================
\section{Resolution: the ZBW Orbit is a Discrete \texorpdfstring{$z$}{z}-Polygon}
\label{sec:resolution}
% ===================================================================

The CPP lattice is discrete by axiom A2 (the
600-cell topology).  Every CP resides on a vertex; propagation
proceeds along edges; angular orbits through closed loops of
nearest-neighbour edges.  There is no continuous angular
coordinate at the lattice level --- ``circular'' paths are
polygonal.

\subsection{Discrete circumference of the ZBW orbit}

The ZBW (zitterbewegung) orbit of a bound qCP charge in a quark
cage is, in CPP, a closed lattice walk through the $z = 12$
nearest neighbours of the cage vertex.  The orbit visits
successive neighbours in a coordinated pattern, returning
after $z$ steps.  The analogy with a continuous circular orbit
of radius $r$ is
\begin{equation}
\underbrace{2\pi r}_{\text{continuum circumference}}
\;\;\longleftrightarrow\;\;
\underbrace{z \cdot \ledge}_{\text{discrete polygon perimeter}},
\label{eq:discrete_correspondence}
\end{equation}
since a $z$-sided polygon of edge length $\ledge$ has perimeter
$z \ledge$.  In the $z \to \infty$ continuum limit with edge
length $\ledge \to 0$ at fixed circumference, the discrete
perimeter approaches $2\pi r$ with $r = \ledge/(2\sin(\pi/z))
\to z\ledge/(2\pi)$; but the discrete lattice does not realise
that limit.

\subsection{What SS-2 Counted}

The SS-2 heuristic \eqref{eq:heuristic} grouped its factors as
\begin{equation}
\frac{M_0}{\ledge}
\;\text{(linear DP energy density)}
\;\times\;
\frac{z\pi}{\phig}
\;\text{(ostensibly orbital factor)},
\end{equation}
with ``$\pi$ from the ZBW circular orbit'' covering an undifferentiated
dimensionless multiplicity.  Given that the 600-cell carries no
continuum angular coordinate at the lattice level, the $\pi$ is
not being used in its strict geometric sense of half-circumference
per radius; it is used as a stand-in constant of the right order of
magnitude.  The relation \eqref{eq:discrete_correspondence}
identifies the correct discrete analogue: any genuine factor of
$\pi$ on the continuum side corresponds to a factor of $z$ (polygon
vertex count) on the discrete side.  Replacing the heuristic $\pi$
by the lattice-native $z$ yields $z \cdot z = z^2$ in place of
$z\pi$.

The replacement is heuristic-to-heuristic at this stage: it promotes
a continuum stand-in to a discrete stand-in that happens to give the
correct numerical value.  A non-heuristic reading of $z^2$ is given
next.

\subsection{What $z^2$ counts}
\label{ssec:zsquared}

A cleaner physical reading of $z^2$, independent of the
heuristic-repair argument, is as the mode multiplicity per
lattice edge in a flux tube that runs along the lattice:
\begin{itemize}[leftmargin=1.5em]
\item \textbf{Along the tube axis.} At each vertex along the tube,
  the chain can connect to $z$ nearest neighbours.  One of those
  is the next vertex along the axis; the remaining $z - 1$
  would be transverse, but each is reciprocally paired across
  the bond, and the effective axial coordination is $z$ in the
  sense of directed-mode labelling $\Tr A^2 / N_V = 12 = z$
  (with $A$ the 600-cell adjacency matrix, $\Tr A^2 = 2E = 1440$,
  $N_V = 120$; see Table~\ref{tab:traces}).
\item \textbf{Around the tube cross-section.} A flux tube of
  radius $\sim \ledge$ on the 600-cell is bounded by the
  second-neighbour shell of the axial vertex chain.  Each vertex
  in that shell contributes a face-mode pair to the tube
  cross-section.  The face-per-vertex ratio on the 600-cell is
  $3F/N_V = 3 \cdot 1200/120 = 30$, of which the $z = 12$
  bounding faces participate (the remaining 18 lie on further
  shells or outside the tube cross-section).
\end{itemize}
Combining the axial and cross-sectional coordination gives
$z \cdot z = z^2$ face-mode couplings per lattice edge of tube
length, each carrying DP energy quantum $M_0$.

This reading \emph{motivates} $z^2$ but does not \emph{derive}
it from the 600-cell trace structure.  The rigorous derivation
is deferred to Section~\ref{sec:open}.

\begin{table}[h]
\centering
\begin{tabular}{@{}lcc@{}}
\toprule
Quantity & Value & Identity \\
\midrule
$N_V$ (vertices) & $120$ & -- \\
$E$ (edges) & $720$ & $N_V z / 2$ \\
$F$ (triangular faces) & $1200$ & $5 N_V / 2$ \\
$C$ (tetrahedral cells) & $600$ & $5 N_V$ \\
$z$ (coordination) & $12$ & $2E/N_V$ \\
$\Tr A^2$ (2 $\times$ edge count) & $1440$ & $2E = N_V z$ \\
$\Tr A^3$ (closed triangles $\times 6$) & $7200$ & $6F$ \\
$\Tr A^2 / N_V$ & $12$ & $z$ \\
$\Tr A^3 / (3 N_V)$ & $20$ & $F \cdot 3 / (3 N_V) \cdot 3 = 20$ \\
\bottomrule
\end{tabular}
\caption{600-cell combinatorial data relevant to the
$z^2$ face-mode counting.  The adjacency-matrix traces
$\Tr A^2 = 2E = 1440$ and $\Tr A^3 = 6F = 7200$ are the
standard inputs to the CPP gauge-coupling formula
$N = \Tr A^2 + \Tr A^3/3 = 1440 + 2400 = 3840$.  The rigorous
derivation of the $z^2$ face-mode multiplicity as a specific
combination of $\Tr A^2$, $\Tr A^3$, and $N_V$ is stated in
Section~\ref{sec:open}.}
\label{tab:traces}
\end{table}

% ===================================================================
\section{Main Conjecture}
\label{sec:main}
% ===================================================================

\begin{conjecture}[String tension from face-mode multiplicity]
\label{conj:main}
The QCD string tension in the CPP framework is
\begin{equation}
\boxed{
\sigma = \frac{M_0 \cdot z^2}{\phig \cdot \ledge}
},
\label{eq:main_conj}
\end{equation}
with $M_0 = m_e z /\phig$, $z = 12$, $\phig$ the golden ratio,
and $\ledge = \lunit/\phig = \hbar c /(\phig \LQCD)$.  Equivalently,
\begin{equation}
\sigma
= \frac{m_e\, z^3}{\phig \, \lunit}
= \frac{m_e\, z^3\, \LQCD}{\phig\, \hbar c}.
\label{eq:main_alt}
\end{equation}
\end{conjecture}

\subsection{Numerical evaluation}

With the values in Section~\ref{sec:primitives},
\begin{align}
\sigma
&= \frac{3.790~\mathrm{MeV} \times 144}{1.618 \times 0.364~\mathrm{fm}} \\
&= \frac{545.8~\mathrm{MeV}}{0.589~\mathrm{fm}} \\
&= 926.5~\MeVfm. \notag
\end{align}

\subsection{Comparison with experiment}

\begin{table}[h]
\centering
\begin{tabular}{@{}lccc@{}}
\toprule
Source & $\sigma$ (MeV/fm) & Basis & Error vs.\ this paper \\
\midrule
SS-2 heuristic, $z\pi/\phig$ & $243$ &
  Continuum ZBW orbit ansatz & $-73.8\%$ \\
This paper, $z^2/\phig$ &
  $\mathbf{926.5}$ &
  Discrete lattice polygon & --- \\
Charmonium Cornell fit &
  $\approx 910$ &
  $c\bar c$ spectrum~\citep{eichten_quigg} & $+1.82\%$ \\
Bottomonium Cornell fit &
  $\approx 900$ &
  $b\bar b$ spectrum~\citep{bali_review} & $+2.94\%$ \\
Lattice QCD (quenched) &
  $\approx 900$--$940$ &
  Static quark-antiquark potential~\citep{bali_review} &
  within range \\
\bottomrule
\end{tabular}
\caption{CPP prediction compared with standard values of
$\sigma$.  The conjectured value $926.5~\MeVfm$ sits within
the lattice-QCD quenched window and within $2$--$3\%$ of the
charmonium and bottomonium Cornell fits.}
\label{tab:comparison}
\end{table}

\begin{remark}[The 1.82\% residual]
\label{rem:residual}
The residual between the conjectured value and the
Cornell fit is consistent with the generic CPP
residual scale of $2$--$4\%$ inherited from the
stereographic projection factor $\phig^{1/z} - 1
\approx 4.1\%$ that also appears in SS-1's derivation of
$\sea$ (SS-1, Table~11.1).  If the conjectured form
\eqref{eq:main_conj} is exact, the full projection correction
applied at the same order should reduce the residual to
the sub-percent level.  A detailed calculation is deferred
to a future revision.
\end{remark}

% ===================================================================
\section{Factorised Physical Interpretation}
\label{sec:factorization}
% ===================================================================

The conjecture \eqref{eq:main_conj} has a natural factorisation
into three physically meaningful pieces:
\begin{equation}
\sigma
\;=\;
\underbrace{\frac{M_0}{\ledge}}_{\text{linear DP density}}
\;\;\times\;\;
\underbrace{z^2}_{\text{face-mode multiplicity}}
\;\;\times\;\;
\underbrace{\eta = \frac{1}{\phig}}_{\text{propagation efficiency}}.
\label{eq:factorisation}
\end{equation}

\begin{description}[style=unboxed,leftmargin=1.5em]
\item[Linear DP density.] A qDP chain carries one DP energy
  quantum $M_0$ per lattice edge of length $\ledge$.  Numerically
  $M_0/\ledge = 3.790/0.364 = 10.41~\MeVfm$.  This is the
  \emph{bare} chain tension in the absence of mode multiplicity.
  SS-2 Section~2 (``Route 1: QCD String Tension'') identified this
  ratio with $\sigma$ directly, obtained $\ledge \sim 0.004$~fm,
  and rejected the route as ``$100\times$ too small because a
  chain contains many DPs per edge length, not one.''  The present
  factorisation identifies the correct many-to-one factor
  explicitly: it is the face-mode multiplicity $z^2 = 144$, which
  is near-exactly the factor of $\sim 90$ (ratio $z^2/\phig$) that
  SS-2's Route~1 was missing.

\item[Face-mode multiplicity.] $z^2 = 144$ is the number of
  face-mode coupling channels per lattice edge of flux tube,
  counting axial $\times$ cross-sectional coordination as argued
  in Section~\ref{ssec:zsquared}.  This replaces the SS-2 factor
  $z\pi \approx 37.70$.  Rigorous derivation deferred
  (Section~\ref{sec:open}).

\item[Propagation efficiency.] $\eta = 1/\phig$ is the standard
  CPP efficiency constant, arising as the 600-cell edge-to-circumradius
  ratio and appearing throughout the gauge-coupling formulas
  $\sin^2\theta_W = 3/(8\phig)$, $\alpha_s = 5/(8\phig)$, and
  the quark mass formula $M_q = m_e(z/\phig)V^{7/3}$~\citep{cpp_sm8,cpp_sm6}.
  Its recurrence here is not ad hoc: whenever a 600-cell mode
  sum is evaluated at the lattice edge scale, the same $\eta$
  projection factor appears.
\end{description}

The factorisation makes explicit that the $z^2$ is the
quantitative content of the conjecture; $M_0/\ledge$ and
$1/\phig$ are already settled within CPP.

% ===================================================================
\section{Cornell-Closure Consistency}
\label{sec:closure}
% ===================================================================

Cornell-closure \eqref{eq:cornell_closure} relates $\sigma$ and
$\rconf$ through $\alpha_s$.  With the conjectured $\sigma =
0.9265~\GeVfm$ and the running coupling $\alpha_s(m_Z) = 0.118$
used in the C14 analysis~\citep{c14,pdg2024}:
\begin{equation}
\rconf
= \sqrt{\frac{\alpha_s\, \hbar c}{\sigma}}
= \sqrt{\frac{0.118 \times 0.1973}{0.9265}}~\mathrm{fm}
= 0.1584~\mathrm{fm}.
\label{eq:rconf}
\end{equation}
The C14 self-consistent solution is $\rconf = 0.161$~fm (at the
target $\sigma = 0.900$~GeV/fm).  The error is $-1.6\%$, of the
same size as the residual on $\sigma$ and consistent with the
generic CPP residual scale.

\begin{remark}[On the choice $\alpha_s = 0.118$ in the closure]
\label{rem:alpha_choice}
The use of $\alpha_s(m_Z) = 0.118$ (rather than a larger hadronic-scale
value) in \eqref{eq:cornell_closure} follows C14 and the standard
quarkonium-spectroscopy convention in which $\alpha_s$ and $\sigma$
are both fitted to the Cornell form with $\alpha_s$ treated as an
effective single-scale coupling.  In CPP this is consistent with
the SS-2 result $\alpha_s(m_Z) = 0.1132$ (derived from
$\alphageom = 1/\sqrt{5}$ running to $m_Z$) and with SM-7's
$\alpha_s = 5/(8\phig) = 0.386$ evaluated at the geometric scale
(not the Cornell-fit scale).  A unified derivation reconciling the
three values (0.113, 0.118, 0.386) with the PSR-saturation running
scheme is OPEN-SS-7 and is cross-referenced
below~(\S\ref{sec:open}).
\end{remark}

\subsection{Simultaneous match on both quantities}

The conjecture \eqref{eq:main_conj} and the closure
\eqref{eq:rconf} match the Cornell fit simultaneously:
\begin{center}
\begin{tabular}{@{}lcc@{}}
\toprule
Quantity & CPP (this paper) & Cornell fit \\
\midrule
$\sigma$ (GeV/fm) & $0.9265$ & $\approx 0.910$ \\
$\rconf$ (fm) & $0.1584$ & $\approx 0.161$ \\
\bottomrule
\end{tabular}
\end{center}
Both emerge at the $\sim 2\%$ level from the single conjectured
formula \eqref{eq:main_conj}, without parameter fitting beyond
the CPP primitives.  The simultaneous match on two independent
quantities is a stronger consistency check than matching
either alone.

% ===================================================================
\section{Physical Mechanism Bridge: CPP Flux Tube \texorpdfstring{$\leftrightarrow$}{<->} QCD Flux Tube}
\label{sec:bridge}
% ===================================================================

Following the template established in SS-3~\S7, we give the
structural (not literal) mapping between the CPP mechanistic
account and the standard QCD picture.

\begin{center}
\begin{tabular}{@{}p{5.8cm}p{5.8cm}@{}}
\toprule
\textbf{QCD (field-theoretic)} & \textbf{CPP (mechanistic)} \\
\midrule
Wilson area law
  $\langle W(C)\rangle \sim e^{-\sigma A}$
  for large rectangular loops of area $A$~\citep{wilson1974} &
  Energy accumulation proportional to enclosed lattice
  area from alternating-stress qDP chains spanning the loop
  interior \\[6pt]
Flux tube of transverse area $\sim 1/\sigma \sim r_{\mathrm{tube}}^2$ &
  Self-collimated qDP chain of transverse extent $\sim \ledge$
  (second-neighbour shell of the axial vertex chain) \\[6pt]
Colour-electric energy density $\frac{1}{2}\vec E_c^2$
  inside the tube &
  Alternating compressive/tensile DP bonds carrying
  pre-stress energy $M_0$ per edge; face-mode multiplicity
  $z^2$ per edge length \\[6pt]
String tension $\sigma$ as phenomenological coefficient
  of the linear Cornell potential &
  Bow-rigidity coefficient set by mode-count
  $\sigma = (M_0/\ledge)(z^2/\phig)$
  from the 600-cell face structure \\[6pt]
String breaking via Schwinger $q\bar q$
  pair production at threshold field strength &
  VP impacts from the DP Sea exceed the bond strength
  at the chain midpoint (central-break dominance
  $\sim 85\%$; see
  \texttt{chain\_fraying\_dynamics.ipynb}) \\[6pt]
Short-distance Coulomb regime
  $V \sim -\alpha_s/r$ for $r \lesssim \rconf$ &
  Sub-threshold regime where self-collimation fails:
  DP Sea cannot nucleate chains fast enough at
  $r \lesssim \rconf$ (PSR saturation; SS-1 Remark~5.3) \\
\bottomrule
\end{tabular}
\end{center}

\begin{remark}[Why the mapping is structural, not literal]
The Wilson area law is a statement about the expectation value
of a path-ordered exponential of gauge fields integrated around
a closed loop.  The CPP account is a deterministic statement
about the energy stored in discrete DP bond configurations on
the lattice interior of the same loop.  Both give the same
dimensional form $\sigma \times A$ for the leading area-law
contribution because both describe the same underlying
structure: a two-dimensional sheet of correlated mode energy
spanning the loop.  The coefficient $\sigma$ in both pictures is
the same quantity, derived here for the first time from the
primitive CPP parameters.
\end{remark}

\begin{remark}[Relation to the bow-rigidity picture of C14]
The bow-rigidity account in C14 and the
\texttt{chain\_fraying\_dynamics} notebook provides the
\emph{shape} of $V(r)$: linear confinement emerges because the
chain bows transversely, storing elastic energy proportional to
extension.  This paper provides the \emph{magnitude}: the
elastic coefficient of that bowing is fixed by the face-mode
multiplicity per lattice edge, not by the bow amplitude
(which is a dimensionless shape parameter in the notebook).  The
two accounts are thus complementary: bow rigidity gives the
$r$-dependence, mode-counting gives the $\sigma$.
\end{remark}

\begin{remark}[The $4+4$ standing-wave extension --- a natural
  sequel, not this paper]
\label{rem:standing_wave}
SS-3 identified the eight gluon degrees of freedom with a
physical basis of four linear DP-chain bond modes plus four
coupled harmonic junction modes on the tetrahedral cage
(the $4 + 4 = 8$ decomposition).  The present paper fixes the
tension $\sigma$ that enters those modes' dynamics.  A natural
physical extension, not pursued quantitatively here, is to
treat the baryon as a closed \emph{standing-wave} configuration
of the $4 + 4$ mode system on the tetrahedral cage.  With
$\sigma$ setting the tension on the linear bonds and the
coupling matrix of SS-3 (with $\det M = 2/\sqrt{3}$) setting
the junction mode structure, the joint standing-wave resonance
condition selects a specific cage edge length $L_{\mathrm{cage}}$:
\begin{equation*}
\{\omega_n(\sigma, \mu_{\mathrm{DP}}, L_{\mathrm{cage}})\}_{n=1}^{8}
\; \longmapsto \;
\text{closed-cage resonance}
\; \longmapsto \;
L_{\mathrm{cage}} = f(\sigma, \mu_{\mathrm{DP}}, \phig).
\end{equation*}
If this resonance condition has a unique stable solution at the
observed $L_{\mathrm{cage}} \sim 0.4$--$0.9~\mathrm{fm}$, then
the proton/neutron radius is determined by the same
mode-counting architecture that determines $\sigma$, without
additional parameters.  The natural scale check
$\sigma \cdot \ledge \approx 0.93 \times 0.364 \approx 340~\mathrm{MeV}$
matches the constituent quark mass scale, which is encouraging
but not conclusive.

This extension is complementary to SS-2's force-balance
derivation of $r_p$ (which used EM repulsion, colour Coulomb,
linear confinement, and a kinetic term, obtaining $r_p = 0.883
~\mathrm{fm}$, $+5.0\%$ vs.\ PDG).  Agreement between the two
routes would be a strong internal consistency check;
disagreement would be diagnostic.  A quantitative development
of the standing-wave proposition is registered as a
conjecture in the Research Frontier and is the natural subject
of a future paper in the Strong Sector series.
\end{remark}

% ===================================================================
\section{Falsifiable Predictions}
\label{sec:predictions}
% ===================================================================

Three predictions follow from combining Conjecture~\ref{conj:main}
with the qualitative mechanism established in the
\texttt{chain\_fraying\_dynamics} and \texttt{zbw\_magnetic\_effects}
notebooks.  The first two are inherited from earlier work; the third
is a specific consequence of the mode-count form of $\sigma$.

\begin{enumerate}[leftmargin=1.8em]

\item \textbf{Central-break dominance in string-breaking events
      ($\sim 85\%$).}
      The differential terminus force $F_{\mathrm{diff}}$ vanishes
      at the midpoint of a qDP chain because the quark-end and
      antiquark-end contributions cancel maximally there.  CPP
      therefore predicts that string-breaking events preferentially
      occur near the chain midpoint, giving daughter mesons of
      similar invariant mass.  Standard QCD (Schwinger mechanism)
      predicts a uniform string-breaking distribution.
      Distinguishable in high-statistics lattice simulations of
      static $q\bar q$ systems at separation
      $r \sim 2\rconf$--$3\rconf$.
      \emph{Source:} \texttt{chain\_fraying\_dynamics.ipynb} Stage~13.

\item \textbf{Helical jet signature in spin-polarised mesons.}
      ZBW Lorentz forces in the chain are perpendicular-dominant
      with a helical phase pattern coherently oriented along the
      chain axis for polarised parent states.  This generates
      a measurable azimuthal asymmetry in jet fragmentation
      products, detectable at LHCb with polarised $B$-meson samples.
      No corresponding prediction in standard QCD.
      \emph{Source:} \texttt{zbw\_magnetic\_effects.ipynb} Stage~16.

\item \textbf{Scaling of $\sigma$ with 600-cell analogues at
      modified coordination.}
      The conjecture \eqref{eq:main_conj} asserts
      $\sigma \propto z^2$.  Should lattice simulations be run on
      $H_4$-adjacent tilings with modified coordination
      $\tilde z \neq 12$ (e.g., the $\{3,3,5\}$ truncations or
      tesseract-based 4D analogues), CPP predicts
      $\tilde \sigma / \sigma = (\tilde z/z)^2 \cdot
      (\tilde \ledge / \ledge)^{-1}$ at the same $\LQCD$
      anchoring, a specific and non-trivial scaling law not
      present in continuum QCD.
      \emph{Source:} Conjecture~\ref{conj:main} above.

\end{enumerate}

% ===================================================================
\section{The Remaining Rigorous Step}
\label{sec:open}
% ===================================================================

The central conjecture of this paper is the identification of the
dimensionless face-mode multiplicity per lattice edge as $z^2$.
The physical motivation in Section~\ref{ssec:zsquared} is
plausibility, not proof.  The rigorous version would derive $z^2$
from the 600-cell combinatorial data in Table~\ref{tab:traces}.

\begin{openstep}[Derive $z^2$ as a specific combination of
600-cell traces]
\label{open:zsquared}
Show, as a theorem about the 600-cell adjacency matrix $A$, that
the face-mode multiplicity per edge appearing in the flux-tube
tension is
\begin{equation}
\mu_{\mathrm{face}}
\;=\;
z^2
\;=\;
\frac{(\Tr A^2)^2}{N_V \cdot \Tr A^2}
\;=\;
\frac{\Tr A^2}{N_V},
\end{equation}
which equals $z$ only at the identity level and must pick up an
additional factor of $z$ from cross-sectional mode multiplicity to
reach $z^2$.  Specifically: (i)~identify the subset of faces of
the 600-cell that bound a flux tube of minimum transverse extent
($z$ faces per axial edge, conjectured); (ii)~show that each
bounding face contributes one face-mode coupling $M_0$ to the
tube tension; (iii)~combine the $z$ axial coordinations with the
$z$ cross-sectional faces to produce $z^2$ total couplings per
edge.  Numerical verification on the explicit 600-cell adjacency
matrix (Appendix~A.2 of SS-1) is expected to confirm the counting
immediately; the \emph{analytic} derivation requires a face-labelling
convention and a choice of tube axis, neither of which is canonical
on the 600-cell.
\end{openstep}

Completing this step elevates Conjecture~\ref{conj:main} to a
Theorem.  Candidate cross-checks once the counting is analytically
fixed:
\begin{itemize}[leftmargin=1.5em]
\item Relation to $\alpha_s = 5/(8\phig) = \eta \Tr A^3 /(3N)$ with
  $N = 3840$: does the same mode-counting give the gauge coupling
  consistent with SS-1~\S5?
\item Compatibility with the SS-3 $4+4$ mode decomposition: the
  eight Gell-Mann generators split as 4 linear bond modes + 4
  harmonic junction modes.  The $z^2 = 144$ face-mode count
  should factorise as a Kronecker product of these 8 modes with
  the 18 cross-sectional directions that are not on the axis,
  summing to $8 \cdot 18 = 144$.  If this factorisation holds, it
  provides the analytic structure.
\end{itemize}

\subsection{Cross-references and dependencies}

This paper's result feeds directly into:
\begin{itemize}[leftmargin=1.5em]
\item OPEN-SS-6 (glueball mass from closed tetrahedral hDP loop):
  once $\sigma$ is fixed, the closed-loop energy $\sim \sigma
  \cdot P$ with $P$ the loop perimeter gives a direct glueball
  mass prediction.
\item OPEN-SS-7 ($\LQCD$ from PSR saturation): closes the
  last independent hadronic-scale parameter once $\sigma$ and
  $\rconf$ are both from CPP primitives.
\item OPEN-SS-10 (nuclear binding energy $V(r)$ from qDP chain
  insertion): the internucleonic force range is set by $\rconf$.
\item OPEN-SS-14 (QCD deconfinement temperature):
  $T_c \sim \sqrt{\sigma/\pi}$ dimensional estimate gives
  $T_c \approx 170$~MeV, to be refined.
\end{itemize}

% ===================================================================
\section{Conclusion}
\label{sec:conclusion}
% ===================================================================

The CPP framework derives the QCD string tension from its
primitive parameters:
\begin{equation}
\sigma = \frac{M_0 \cdot z^2}{\phig \cdot \ledge}
= \frac{m_e \cdot z^3 \cdot \LQCD}{\phig \cdot \hbar c}
= 926.5~\MeVfm,
\end{equation}
in agreement with the charmonium Cornell fit to $+1.82\%$ and
with the self-consistent Cornell-closure confinement radius
$\rconf = 0.158~\mathrm{fm}$ to $-1.6\%$.  The result corrects
SS-2's heuristic $\sigma = M_0 z\pi/(\phig \ledge) = 243~\MeVfm$
by identifying the ``ZBW circular orbit'' as a discrete $z$-polygon
on the lattice rather than a continuous circle.  The factor
$z/\pi \approx 3.820$ by which SS-2 undershot the empirical value
is physically identified as the discrete-to-continuum correction
on the orbit.

Physically, $\sigma$ factorises as a linear DP energy density
$(M_0/\ledge \approx 10~\MeVfm)$ enhanced by the face-mode
multiplicity per edge $(z^2 = 144)$ and moderated by the
standard 600-cell propagation efficiency $(\eta = 1/\phig)$.

The formula is registered as a Conjecture pending the rigorous
derivation of the $z^2$ multiplicity from the 600-cell trace
structure (Section~\ref{sec:open}).  Numerical agreement at the
$2\%$ level on two independent Cornell quantities provides strong
evidence that the conjecture is correct; the outstanding step is
an analytic proof of the face-labelling count, not a revision of
the formula.

\subsection{Problem Status After This Paper}

\begin{itemize}[leftmargin=1.5em]
\item \textbf{OPEN-SS-5 (string tension $\sigma$ from $\sea$):}
  OPEN $\to$ \textbf{CONJ-SS-5} (Conjecture~\ref{conj:main}),
  subject to the open step on $z^2$ multiplicity
  (OPEN-SS-15, registered by this paper).
  Per CPP nomenclature, the problem number~5 is retained through
  the progression OPEN $\to$ CONJ $\to$ (eventually) THEO.
\item \textbf{CONJ-SS-2-1} (SS-2's heuristic $\sigma = M_0 z\pi/(\phig\ledge)
  = 243~\MeVfm$) $\to$ \textbf{FALS-SS-2-1} (superseded by
  Conjecture~\ref{conj:main}).  The $z/\pi$ factor is
  structurally identified, not empirically discarded; SS-2's
  proton radius calculation inherits a quantitatively correct
  string tension, though $\varepsilon$ and $r_p$ will need
  re-derivation with the new $\sigma$ in the next SS-2 revision.
\item \textbf{OPEN-SS-15 (rigorous $z^2$ face-mode multiplicity
  from 600-cell traces):}
  NEW (registered in Research\_Frontier.md \S1 by this paper).
  Completing this step elevates CONJ-SS-5 to THEO-SS-[next].
\end{itemize}

% ===================================================================
\section*{Acknowledgements}
% ===================================================================

Thomas Lee Abshier, ND, conceived the CPP framework and the
self-consistent closure mechanism in the \texttt{chain\_fraying}
and \texttt{zbw\_magnetic} notebooks that provided the qualitative
basis for this paper.  The SS-2 heuristic that this paper
supersedes was proposed by TLA and Claude Opus in March 2026 and
explicitly flagged there as non-rigorous.  The identification of
the $z/\pi$ correction factor and the discrete-polygon physical
reading emerged from a session between TLA and Claude Opus on
16 April 2026.

This paper is a v0.1 draft and has not yet been through the
Copilot $\to$ Grok $\to$ Sonnet review cycle.  Revisions are
expected.

% ===================================================================
\section{Appendix A: Numerical Verification}
\label{sec:appendix}
% ===================================================================

\subsection{Primitive values}

The derivation uses:
\begin{align*}
m_e       &= 0.51100~\mathrm{MeV} \\
z         &= 12 \\
\phig     &= 1.61803 \\
\LQCD     &= 335~\mathrm{MeV} \quad \text{(SS-2)} \\
\hbar c   &= 0.19733~\mathrm{GeV\cdot fm}
\end{align*}
Derived:
\begin{align*}
\alphageom
  &= 3(11 + 5\sqrt{5})\sqrt{5+\sqrt{5}}/320 = 0.55936 \\
\kSM      &= \alphageom/(z \phig^2) = 0.017805 \\
\sea      &= 10 \kSM = 0.17805 \\
M_0       &= m_e z/\phig = 3.790~\mathrm{MeV} \\
\lunit    &= \hbar c/\LQCD = 0.5890~\mathrm{fm} \\
\ledge    &= \lunit/\phig = 0.3641~\mathrm{fm}
\end{align*}

\subsection{Central computation}

\begin{align*}
\sigma
&= \frac{M_0 z^2}{\phig \ledge} \\
&= \frac{3.790 \times 144}{1.61803 \times 0.3641}~\mathrm{MeV/fm} \\
&= \frac{545.76}{0.5890}~\mathrm{MeV/fm} \\
&= 926.58~\MeVfm.
\end{align*}

Equivalent form via $\LQCD$:
\begin{align*}
\sigma
&= \frac{m_e z^3 \LQCD}{\phig \hbar c} \\
&= \frac{0.51100~\mathrm{MeV} \times 1728 \times 0.335~\mathrm{GeV}}
        {1.61803 \times 0.19733~\mathrm{GeV\cdot fm}} \\
&= \frac{0.29586~\mathrm{GeV}^2}
        {0.31923~\mathrm{GeV \cdot fm}} \\
&= 0.9268~\GeVfm.
\end{align*}
The two forms agree to the displayed precision.

\subsection{Cornell-closure confinement radius}

\begin{align*}
\rconf
&= \sqrt{\alpha_s \hbar c / \sigma} \\
&= \sqrt{\frac{0.118 \times 0.1973}{0.9266}}~\mathrm{fm} \\
&= \sqrt{0.02513}~\mathrm{fm} \\
&= 0.1585~\mathrm{fm}.
\end{align*}
C14 target: $\rconf = 0.161$~fm.  Error: $-1.55\%$.

\subsection{Error tabulation}

\begin{table}[h]
\centering
\begin{tabular}{@{}lccc@{}}
\toprule
Quantity & CPP & Empirical/Cornell & Error \\
\midrule
$\sigma$ (MeV/fm) & $926.58$ & $\approx 910$ & $+1.82\%$ \\
$\rconf$ (fm) & $0.1585$ & $\approx 0.161$ & $-1.55\%$ \\
Product $\sigma \cdot \rconf^2$ (MeV$\cdot$fm) & $23.27$ & $23.59$ & $-1.36\%$ \\
\bottomrule
\end{tabular}
\caption{Consistency tests at the $\sim 2\%$ level.  The product
$\sigma \cdot \rconf^2$ equals $\alpha_s \hbar c = 0.118 \times
197.3 = 23.28$~MeV$\cdot$fm, reflecting the Cornell closure.}
\label{tab:errors}
\end{table}

% ===================================================================
% Bibliography
% ===================================================================

\begin{thebibliography}{99}

\bibitem[Abshier and Grok (2026a)]{cpp_ss1}
Abshier, T.~L. and Grok. (2026).
\emph{SS-1: The Strong Sector from the 600-Cell Lattice}.
CPP Strong Sector Series, v2, Hyperphysics Institute.
\url{https://github.com/Hyperphysics-Institute/CPP}.

\bibitem[Abshier and Opus (2026a)]{cpp_ss2_lattice}
Abshier, T.~L. and Claude Opus. (2026).
\emph{SS-2: Lattice-Scale Grounding and Nucleon Structure}.
CPP Strong Sector Series, v1.0, Hyperphysics Institute.

\bibitem[Abshier and Opus (2026b)]{cpp_ss3}
Abshier, T.~L. and Claude Opus. (2026).
\emph{SS-3: Uniqueness of SU(3) from the Tetrahedral Cage}.
CPP Strong Sector Series, v1.3, Hyperphysics Institute.

\bibitem[Abshier and Opus (2026c)]{cpp_sm8}
Abshier, T.~L. and Claude Opus. (2026).
\emph{SM-8: Quark Generation Structure from 600-Cell Distance Shells}.
CPP Standard Model Emergence Series, v4.1, Hyperphysics Institute.

\bibitem[Abshier and Opus (2026d)]{cpp_sm6}
Abshier, T.~L. and Claude Opus. (2026).
\emph{SM-6: Charged Lepton Mass Spectrum from 600-Cell Lattice Geometry}.
CPP Standard Model Emergence Series, v3, Hyperphysics Institute.

\bibitem[Abshier and Grok (2026b)]{c14}
Abshier, T.~L. and Grok. (2026).
\emph{C14: Cornell Potential Closure from qDP Chain Dynamics}.
CPP Companion Paper, Hyperphysics Institute.

\bibitem[Eichten and Quigg (2008)]{eichten_quigg}
Eichten, E. and Quigg, C. (2008).
Mesons with beauty and charm: Spectroscopy.
\emph{Physical Review D}, 49, 5845--5856
(and updates via PDG).

\bibitem[Bali (2001)]{bali_review}
Bali, G.~S. (2001).
QCD forces and heavy quark bound states.
\emph{Physics Reports}, 343, 1--136.

\bibitem[Wilson (1974)]{wilson1974}
Wilson, K.~G. (1974).
Confinement of quarks.
\emph{Physical Review D}, 10, 2445--2459.

\bibitem[Particle Data Group (2024)]{pdg2024}
Particle Data Group. (2024).
Review of Particle Physics.
\emph{Progress of Theoretical and Experimental Physics}, 083C01.

\end{thebibliography}

% BEGIN GENERATED IDENTIFIER APPENDIX -- do not edit by hand
\clearpage
\section*{Appendix: Programme identifiers used in this paper}
\addcontentsline{toc}{section}{Appendix: Programme identifiers}

\noindent This programme tracks its unresolved questions under short identifiers so that a claim and its open dependencies can be cited precisely. The identifiers appearing in this paper are glossed below; no external document is needed to read them.

\begin{description}
  \item[\texttt{OPEN-SS-10}] \hfill \\ Derive nucleon–nucleon potential V(r) from qDP chain insertion dynamics.
  \item[\texttt{OPEN-SS-14}] \hfill \\ Derive T\_c ≈ 155 ± 10 MeV from sea\_strength and 600-cell geometry.
  \item[\texttt{OPEN-SS-15}] \hfill \\ The open step on the z-squared multiplicity in the string-tension derivation.
  \item[\texttt{OPEN-SS-5}] \hfill \\ Derive σ ≈ 0.9 GeV/fm from sea\_strength and 600-cell geometry without calibration.
  \item[\texttt{OPEN-SS-6}] \hfill \\ Compute lightest scalar glueball mass using f\_geom formula applied to closed hDP loop.
  \item[\texttt{OPEN-SS-7}] \hfill \\ Derive Λ\_QCD ≈ 0.218 GeV from PSR saturation without PDG input.
\end{description}
% END GENERATED IDENTIFIER APPENDIX

\end{document}
